\section{Error analysis}
\label{sec:order}

The analysis has three parts. Lemma~\ref{lem:exact-stage-identity} shows that the stage system of
Eq.~\eqref{eq:g6fullva-expanded-stage} is exactly three-stage Gauss collocation of the reduced
joint-coordinate equations of motion, written in absolute coordinates; the stage residual vanishes
identically at the lifted reduced Gauss stage and no stage-residual perturbation term appears.
The two perturbations the implementation does introduce, the endpoint closure and the inexact
Newton solve, are bounded by $O(h^7)$ (Lemmas~\ref{lem:endpoint-closure}
and~\ref{lem:inexact-newton}), and the local defect is transferred to the grid by a discrete
Gronwall argument (Lemma~\ref{lem:local-global-transfer}). The hypotheses are collected in
Assumption~\ref{ass:regularity}; Lemmas~\ref{lem:p2-from-p1} and~\ref{lem:newton-envelope} show that
the second and third of them follow from the first for small steps. The algebraic identities and the
perturbation lemmas are machine-checked in Lean~4/Mathlib (\ref{app:lean}); Butcher's
order theorem for the Gauss tableau and the classical Gauss local-error estimate
\cite{hairer1993solving} are used as cited results.

\subsection{Reduced chart, lift and hypotheses}

Let $y=(s_0,\theta_0,s_1,\theta_1,\dot s_0,\dot\theta_0,\dot s_1,\dot\theta_1)$ collect the joint
coordinates of Eq.~\eqref{eq:joint-coordinates} and their rates. On the smooth constrained branch
the equations of motion reduce to
\begin{equation}
  \dot y=f(y,t),
  \label{eq:reduced-ode}
\end{equation}
where $f$ returns the rates and the joint accelerations obtained from the constrained Newton--Euler
equations. Let $\varphi_h$ denote the flow of Eq.~\eqref{eq:reduced-ode} over one step,
$\mathcal L(y,t)=(q,v,a,\lambda)$ the smooth lift to absolute coordinates given by forward
kinematics together with the velocities, accelerations and multipliers of the constrained dynamics,
and $Y_1,Y_2,Y_3$ the Gauss collocation stages of Eq.~\eqref{eq:reduced-ode} from $y_n$,
\begin{equation}
  Y_i=y_n+h\sum_{k=1}^3A_{ik}f(Y_k,t_k),\qquad i=1,2,3 .
  \label{eq:reduced-gauss-stages}
\end{equation}
The lifted reduced Gauss stage is $Z_G=(\mathcal L(Y_1),\mathcal L(Y_2),\mathcal L(Y_3))$. We write
$F_h(\cdot;x_n)$ for the implemented $132$-row residual, $\mathcal E_h$ for the endpoint
reconstruction of Eqs.~\eqref{eq:end-trans}--\eqref{eq:end-rot}, $\mathcal C_h$ for the endpoint
closure, $\mathcal P_h=\mathcal C_h\circ\mathcal E_h$, and $\Psi_h^G=\mathcal E_h(Z_G)$ for the
reconstructed Gauss endpoint. Norms are fixed finite-dimensional norms on the reduced chart, the
stage space and the endpoint space; chart-to-configuration norm equivalences on the compact tube are
absorbed into the constants.

\begin{assumption}[Regularity and perturbation scale]
\label{ass:regularity}
Fix a final time $T$ and a compact reduced-state tube $K$ containing the exact trajectory of
Eq.~\eqref{eq:reduced-ode} at distance at least $d_K>0$ from $\partial K$. The following hold
uniformly for $y\in K$, $0\le t\le T$ and $0<h\le h_0$, with constants independent of $h$ and of
the number of steps.
\begin{enumerate}
  \item[(H1)] (Smooth branch.) The constraint Jacobian $\Phi_q$ has constant full row rank, the
  constrained mass matrix is nonsingular on the tangent space, $f$ is $C^7$, the lift $\mathcal L$
  is smooth on the tube, and the joint coordinates complete the constraints to a chart, i.e.\ the
  map $q\mapsto(\Phi(q),s_0(q),\theta_0(q),s_1(q),\theta_1(q))$ has an invertible Jacobian.
  \item[(H2)] (Local inverses and stability.) The stage Jacobian $J_h=D_ZF_h(Z_G;x_n)$ is
  invertible with $\|J_h^{-1}\|\le M$ and $\|D_Z^2F_h\|\le L$ on a ball $B_\rho(Z_G)$ with
  $ML\rho\le\tfrac12$. The closure subsystem $E$ of Lemma~\ref{lem:endpoint-closure} is $C^2$ on a
  ball $B_{r_E}(z_\ast)$ around a same-branch closed anchor $z_\ast$ with $E(z_\ast)=0$;
  $DE(z_\ast)$ has a right inverse $B_\ast$ with $\|B_\ast\|\le M_{E}$,
  $\|DE(z)-DE(z_\ast)\|\le L_E\|z-z_\ast\|$ on that ball, $M_{E}L_Er_E\le\tfrac12$, the unclosed
  endpoint $z_u=\Psi_h^G$ lies in $B_{r_E/2}(z_\ast)$, and its raw closure defect satisfies
  $\|E(z_u)\|\le C_{E,\mathrm{raw}}h^7$. The one-step map satisfies
  \begin{equation}
  \label{eq:assumption-stability-scale}
    \|\Psi_h(y)-\Psi_h(\bar y)\|\le (1+C_sh)\|y-\bar y\|,\qquad y,\bar y\in K .
  \end{equation}
  \item[(H3)] (Solver envelope.) The computed Newton iterate $\widetilde Z$ lies in the ball
  $B_\rho(Z_G)$ of (H2) and its residual satisfies $\|F_h(\widetilde Z;x_n)\|\le\eta_h\le c_\eta h^7$.
  The output map satisfies $\|D_Z\mathcal P_h\|\le M_N$ on that ball.
\end{enumerate}
\end{assumption}

(H1) is the compact-tube regularity hypothesis of classical collocation theory. (H2) and (H3) are
the interfaces through which the implementation enters; Lemmas~\ref{lem:p2-from-p1}
and~\ref{lem:newton-envelope} derive them from (H1) for $h\le h_0$, and they are kept as separate
items only because the theorem is stated for any mechanism whose stage system has the form of
Eq.~\eqref{eq:g6fullva-expanded-stage}. (In the accompanying repository the three hypotheses are
called P1, P2 and P6.) The theorem also uses the classical Gauss local-error estimate: for
$f\in C^7$ on the compact tube there is $C_{G,0}$ with
\begin{equation}
  \|u(t_n+h)-\varphi_h(y_n)\|\le C_{G,0}h^7,
  \label{eq:gauss-local-truncation}
\end{equation}
where $u$ is the degree-three collocation polynomial through Eq.~\eqref{eq:reduced-gauss-stages}
\cite{hairer1993solving}. The tableau used in the implementation satisfies Butcher's simplifying
assumptions $B(6)$, $C(3)$, $D(3)$ and fails $B(7)$, so its order is exactly six; these algebraic
conditions are machine-checked for the coefficients hard-coded in the implementation.

\subsection{Exact stage identity}

\begin{lemma}[Exact stage identity]
\label{lem:exact-stage-identity}
Under (H1), the implemented $132$-row residual satisfies
\begin{equation}
  F_h(Z_G;x_n)=0 .
  \label{eq:exact-stage-identity}
\end{equation}
Conversely, on the branch where $n_j\cdot a_j(q_i)\neq0$, every stage vector that satisfies the
$96$ non-dynamic rows of Eq.~\eqref{eq:g6fullva-expanded-stage} lies on the constraint manifold
with consistent velocities and accelerations and has joint coordinates satisfying
Eq.~\eqref{eq:reduced-gauss-stages}. Hence the stage solve is three-stage Gauss collocation of
Eq.~\eqref{eq:reduced-ode} in the joint-coordinate chart, written in absolute coordinates.
\end{lemma}

\begin{proof}
The $72$ constraint rows vanish on $Z_G$ because each $\mathcal L(Y_k)$ lies on the constraint
manifold with consistent velocities and accelerations. The $24$ joint-coordinate rows are, by
Eq.~\eqref{eq:joint-coordinates}, the reduced collocation equations~\eqref{eq:reduced-gauss-stages}
read component-wise, because $s_j,\theta_j,\dot s_j,\dot\theta_j$ are coordinates of $y$ and
$\ddot s_j,\ddot\theta_j$ are components of $f$; the implemented projection onto $a_j(q_i)$
multiplies each translational row by $n_j\cdot a_j(q_i)$. The $36$ Newton--Euler rows vanish because
$\mathcal L$ supplies the stage accelerations and multipliers satisfying force and moment balance,
and the implemented rows are exactly these balance defects with the implemented wrench sign
conventions. For the converse, the constraint rows force $d_j,\dot d_j,\ddot d_j\parallel n_j$, so
each translational collocation row factors as
$\bigl(s_j(Z_i)-s_j(x_n)-h\sum_kA_{ik}\dot s_j(Z_k)\bigr)\,(n_j\cdot a_j(q_i))$, and the twist rows
are already scalar. Both directions are verified row by row for a transcription of the implemented
residual in the Lean development.
\end{proof}

The lemma replaces any stage-residual perturbation estimate: the residual at the lifted Gauss
stage is not small, it is zero. What remains to be controlled is which root Newton selects, how the
endpoint is reconstructed and closed, and how inexact the Newton solve is. The lemma is stated for
the open chain, where the joint coordinates form a global chart; for closed loops the same argument
applies on any local chart of independent coordinates. The rows of
Eq.~\eqref{eq:g6fullva-expanded-stage} are transcribed from the single implemented residual
function, whose automatic derivative is the Newton Jacobian, so there is one residual and Jacobian
path; the Lean development transcribes the same rows, and the binding between transcription and
source is a static check, not a theorem. Numerically, the converged stage vectors of the chain runs
read in the joint-coordinate chart have reduced collocation defects and out-of-axis components of
$d_j,\dot d_j,\ddot d_j$ at roundoff (Section~\ref{sec:numerical}).

\begin{lemma}[Branch selection]
\label{lem:stage-residual-defect}
Under (H2), $Z_G$ is the unique zero of $F_h(\cdot;x_n)$ in $B_\rho(Z_G)$. In particular the exact
stage root on that branch equals $Z_G$, and $\mathcal E_h(Z_G)=\Psi_h^G$.
\end{lemma}

\begin{proof}
For $Z\in B_\rho(Z_G)$ the mean value inequality gives
$\|F_h(Z)-F_h(Z_G)-J_h(Z-Z_G)\|\le\tfrac{L\rho}{2}\|Z-Z_G\|$. If $F_h(Z)=0$, then
Eq.~\eqref{eq:exact-stage-identity} gives
$\|Z-Z_G\|\le M\|F_h(Z)-F_h(Z_G)-J_h(Z-Z_G)\|\le\tfrac{ML\rho}{2}\|Z-Z_G\|\le\tfrac14\|Z-Z_G\|$,
so $Z=Z_G$.
\end{proof}

\begin{lemma}[(H2) from (H1) for small steps]
\label{lem:p2-from-p1}
Assume (H1). Then there is $h_0>0$ such that for $0<h\le h_0$: (i) the stage Jacobian
$J_h=D_ZF_h(Z_G)$ is invertible with
\begin{equation}
  \|J_h^{-1}\|\le 2\|J_0^{-1}\|,\qquad J_0=\lim_{h\to0}J_h,\qquad \|J_h-J_0\|\le C_Jh ,
  \label{eq:p2-perturbation-bound}
\end{equation}
(ii) the closure Jacobian is invertible, so $B_\ast$ exists with $\|B_\ast\|$ bounded uniformly,
(iii) the raw closure defect satisfies $\|E(z_u)\|\le C_{E,\mathrm{raw}}h^7$, and (iv) the one-step
map satisfies Eq.~\eqref{eq:assumption-stability-scale}. Hence all of (H2) follows from (H1) for
small steps.
\end{lemma}

\begin{proof}
At $h=0$ the collocation rows reduce to $s_j(Z_i)-s_j(x_n)$, $\theta_j(Z_i)-\theta_j(x_n)$ and
their rate analogues, so $J_0$ is block diagonal over the stages and block triangular within a
stage: the position block is the Jacobian of $q\mapsto(\Phi,s,\theta)$, the velocity block is the
same matrix acting on $(v,\omega)$ because on the constraint manifold the rate rows are the time
derivatives of the coordinate rows, and the acceleration block is the constrained Newton--Euler
system $\bigl[\begin{smallmatrix}M&\Phi_q^T\\ \Phi_q&0\end{smallmatrix}\bigr]$ acting on
$(a,\alpha,\lambda)$. Each block is invertible under (H1), so $J_0$ is invertible with an
$h$-independent inverse bound. The only $h$-dependence of $F_h$ is the factor $h\sum_kA_{ik}$ in
the collocation rows, so $\|J_h-J_0\|\le C_Jh$ on the compact tube, and for
$2\|J_0^{-1}\|C_Jh\le1$ Eq.~\eqref{eq:p2-perturbation-bound} follows from the Neumann series. The
endpoint closure is a square KKT system in the endpoint velocities whose matrix is
$\bigl[\begin{smallmatrix}I&\Phi_q^T\\ \Phi_q&0\end{smallmatrix}\bigr]$, invertible by full row rank
of $\Phi_q$; its raw defect is the velocity-level constraint value at the Gauss-weight endpoint,
which is the quadrature error of Lemma~\ref{lem:gauss-endpoint-defect} and hence $O(h^7)$; the
implicit function theorem gives a same-branch closed anchor $z_\ast$ within $O(h^7)$ of $z_u$, so
$z_u\in B_{r_E/2}(z_\ast)$ for small $h$. Finally, by Lemma~\ref{lem:exact-stage-identity} the
one-step map read in the reduced chart is the reduced Gauss step, whose derivative is $I+O(h)$ by
the implicit function theorem for the stage equations, composed with the smooth endpoint
reconstruction and an $O(h^7)$ closure correction; this gives the Lipschitz scale $1+C_sh$.
\end{proof}

The threshold $h_0$ of the Neumann argument is not sharp. On the chain
(Table~\ref{tab:p2-constants-check} in \ref{app:diagnostics}) the sufficient condition
$\|J_0^{-1}\|\,\|J_h-J_0\|\le\tfrac12$ would require $h\le2\times10^{-5}$ in the Euclidean norm on
the stage layout and $h\le8\times10^{-4}$ in the endpoint-linearized norm $\|J_0^{-1}\,\cdot\,\|$,
whereas the conclusion $\|J_h^{-1}\|\le2\|J_0^{-1}\|$ holds on every grid used with ratio at most
$1.55$. The gap comes from the curvature of the friction law: the largest entries of $(J_h-J_0)/h$
sit in the Newton--Euler rows against the stage velocities, where the second derivative of the
Brown--McPhee law with $v_s=0.5$ enters, and the same mechanism without joint friction has $C_J$
smaller by a factor of about fifteen and $h_0\approx1.4\times10^{-2}$. The friction law, not the
collocation structure, sets the small-step threshold; Section~\ref{sec:numerical} shows the same
dependence in the observed orders when $v_s$ is decreased.

\subsection{Endpoint defect, closure and inexact Newton}

\begin{lemma}[Gauss endpoint defect]
\label{lem:gauss-endpoint-defect}
Under (H1) there is a constant $C_G$, uniform on the compact tube, such that
\begin{equation}
  \|\Psi_h^G(y)-\varphi_h(y)\|\le C_Gh^7 .
  \label{eq:gauss-endpoint-defect}
\end{equation}
\end{lemma}

\begin{proof}
Let $u$ be the collocation polynomial of Eq.~\eqref{eq:reduced-gauss-stages} and $g$ any
component of the lifted state along $u$, so that $g(t)=\rho(u(t))$ for a smooth chart function
$\rho$ and $g'(t_k)$ equals the lifted stage velocity at $Y_k$; for the rotational component, $g$
is the Lie-algebra increment $\eta(t)$ with $\eta'=J_r^{-1}(\eta)\omega$ along the lifted
trajectory. The reconstruction Eqs.~\eqref{eq:end-trans}--\eqref{eq:end-rot} replaces
$g(t_n+h)-g(t_n)=\int_{t_n}^{t_n+h}g'(t)\,dt$ by the Gauss quadrature $h\sum_ib_ig'(t_i)$. Because
the Gauss weights integrate polynomials of degree at most five exactly, a Peano-kernel argument
gives
\begin{equation}
  \Bigl|\int_{t_n}^{t_n+h}g'(t)\,dt-h\sum_{i=1}^3b_ig'(t_i)\Bigr|
  \le\frac{h^7}{60}\max_{[t_n,t_n+h]}|g^{(7)}| ,
  \label{eq:gauss-quadrature-defect}
\end{equation}
with $g^{(7)}$ bounded uniformly on the tube by (H1) and the smoothness of $u$. Hence each
reconstructed endpoint component equals the lifted endpoint of $u$ up to $O(h^7)$, and $\mathrm{Exp}$
is smooth. Combining with Eq.~\eqref{eq:gauss-local-truncation} and the chart norm equivalence gives
Eq.~\eqref{eq:gauss-endpoint-defect}. The bound Eq.~\eqref{eq:gauss-quadrature-defect} is
machine-checked.
\end{proof}

\begin{lemma}[Endpoint-closure perturbation]
\label{lem:endpoint-closure}
Let $E(z)=0$ be the velocity-level KKT closure subsystem solved by $\mathcal C_h$, and let
$z_u=\Psi_h^G$ be the unclosed endpoint. Under (H2) the right-inverse Newton iteration
$z_{l+1}=z_l-B_\ast E(z_l)$, $z_0=z_u$, converges to a closed endpoint $z_E=\mathcal C_h(z_u)$
with $E(z_E)=0$, and any closed endpoint of the form $z_u+B_\ast w$ in $B_{r_E}(z_\ast)$ satisfies
\begin{equation}
  \|z_E-z_u\|\le 2M_{E}\|E(z_u)\|\le C_Eh^7,\qquad C_E=2M_{E}C_{E,\mathrm{raw}} .
  \label{eq:endpoint-correction-bound}
\end{equation}
\end{lemma}

\begin{proof}
Write $G(w)=E(z_u+B_\ast w)$. Since $DE(z_\ast)B_\ast=I$ and
$\|DE(z)-DE(z_\ast)\|\le L_E\|z-z_\ast\|\le L_Er_E$ on the ball, the mean value inequality gives
$\|G(w_1)-G(w_2)-(w_1-w_2)\|\le M_{E}L_Er_E\|w_1-w_2\|\le\tfrac12\|w_1-w_2\|$ whenever both
arguments stay in the ball. The map $w\mapsto w-G(w)$ is therefore a $\tfrac12$-contraction of the
ball of radius $2\|E(z_u)\|$ into itself, which stays inside $B_{r_E}(z_\ast)$ for $h\le h_0$; its
fixed point $w_\ast$ gives $z_E=z_u+B_\ast w_\ast$ with $E(z_E)=0$. For any root of this form, the
same inequality with $w_2=0$ yields $\|w\|\le2\|E(z_u)\|$, hence
$\|z_E-z_u\|=\|B_\ast w\|\le2M_{E}\|E(z_u)\|$. The raw-defect bound of (H2) completes
Eq.~\eqref{eq:endpoint-correction-bound}.
\end{proof}

\begin{lemma}[Inexact Newton]
\label{lem:inexact-newton}
Under (H2) and (H3),
\begin{equation}
\label{eq:inexact-newton-stage-output-bounds}
\begin{aligned}
  \|\widetilde Z-Z_G\|&\le 2M\eta_h,\\
  \|\mathcal P_h(\widetilde Z)-\mathcal P_h(Z_G)\|&\le C_N\eta_h\le C_Nc_\eta h^7,\qquad C_N=2MM_N .
\end{aligned}
\end{equation}
\end{lemma}

\begin{proof}
With Eq.~\eqref{eq:exact-stage-identity} and the mean value inequality of
Lemma~\ref{lem:stage-residual-defect},
$\|\widetilde Z-Z_G\|\le M\bigl(\|F_h(\widetilde Z)\|+\tfrac{L\rho}{2}\|\widetilde Z-Z_G\|\bigr)$,
so $\|\widetilde Z-Z_G\|\le 2M\eta_h$. The output bound follows from $\|D_Z\mathcal P_h\|\le M_N$ on
the ball and the envelope $\eta_h\le c_\eta h^7$.
\end{proof}

\begin{lemma}[Newton reaches the solver envelope]
\label{lem:newton-envelope}
Under (H1) and (H2) there is $h_0>0$ such that for $0<h\le h_0$ the lifted endpoint predictor
$Z^{(0)}_\ast=(\mathcal L(x_n),\mathcal L(x_n),\mathcal L(x_n))$ lies in $B_{\rho/2}(Z_G)$; for
every predictor $Z^{(0)}\in B_{\rho/2}(Z_G)$ the simplified Newton iteration
$Z^{(k+1)}=Z^{(k)}-J_h^{-1}F_h(Z^{(k)})$ stays in $B_\rho(Z_G)$ and
\begin{equation}
  \|Z^{(k)}-Z_G\|\le 2^{-k}\|Z^{(0)}-Z_G\|,\qquad
  \|F_h(Z^{(k)})\|\le L_F\,2^{-k}\|Z^{(0)}-Z_G\| ,
  \label{eq:newton-envelope-decay}
\end{equation}
where $L_F$ is a Lipschitz constant of $F_h$ on the ball; the same decay holds for the full Newton
iteration of Eq.~\eqref{eq:newton}. Consequently the stopping rule $\|F_h(Z^{(k)})\|\le c_\eta h^7$
is met after $k=O(\log(1/h))$ iterations, and (H3) holds by construction whenever the stage solve
starts in $B_{\rho/2}(Z_G)$.
\end{lemma}

\begin{proof}
The stages of $Z_G$ are $\mathcal L(Y_k)$ with $\|Y_k-y_n\|=O(h)$ on the compact tube, so
$\|Z^{(0)}_\ast-Z_G\|=O(h)$ and $Z^{(0)}_\ast\in B_{\rho/2}(Z_G)$ for small $h$. The linearization
inequality of (H2) makes the simplified Newton map a $\tfrac12$-contraction of $B_\rho(Z_G)$ into
itself with fixed point $Z_G$ (Lemma~\ref{lem:stage-residual-defect}), which gives the first bound
in Eq.~\eqref{eq:newton-envelope-decay}; the second follows from $F_h(Z_G)=0$ and the Lipschitz
bound. For the full Newton iteration, (H2) gives $\|D_ZF_h(Z)-J_h\|\le L\|Z-Z_G\|\le L\rho$ on the
ball, so the perturbation bound of Lemma~\ref{lem:p2-from-p1} yields $\|D_ZF_h(Z)^{-1}\|\le 2M$,
and the standard estimate $\|Z^{(k+1)}-Z_G\|\le ML\|Z^{(k)}-Z_G\|^2\le\tfrac12\|Z^{(k)}-Z_G\|$
gives the same decay. The simplified-Newton decay estimate is machine-checked.
\end{proof}

The predictor of the implementation (Section~\ref{sec:method}) is not $Z^{(0)}_\ast$: it sets the
angular-velocity and angular-acceleration guesses to zero, so its distance to $Z_G$ is of the
order of $\|\omega_n\|+\|\alpha_n\|$ rather than $O(h)$ (about $50$ on the chain,
Table~\ref{tab:predictor-check}), and the first full Newton step from it expands the distance by
$12$ to $18\%$ before the quadratic phase sets in. The lemma therefore covers the implemented run
through the ball membership of (H3), which is confirmed a posteriori on every grid used: the
converged roots are the Gauss roots to roundoff. Replacing the zero angular guesses by the endpoint
values would make the implemented predictor an $O(h)$ predictor and bring the run under the lemma
directly.

The experiments of Section~\ref{sec:numerical} use the fixed tolerance $\tau_N=10^{-11}$ rather
than the $h^7$-scaled rule. On the chain sweep the residual of the accepted iterate is between
$10^{-14}$ and $8\times10^{-12}$ at every step and step size (Table~\ref{tab:e7}); relative to
$h^7$ this is $c_\eta\le10^{-3}$ for $h\ge0.05$ but $c_\eta\approx10^{6}$ at $h=0.003125$, so
the fixed-tolerance runs satisfy (H3) only with a large constant on the finest grids, and the
inexact-Newton term $C_Nc_\eta h^7$ of the theorem is then a pessimistic bound: the observed errors
keep decreasing as $h^6$ down to $10^{-13}$, three orders of magnitude below $2M\|F_h(\widetilde Z)\|$.
One additional Newton iteration per step would bring the accepted residual to roundoff at every
step size; the $h^7$-scaled rule with $c_\eta=10^4$ has been exercised on the chain and gives the
same observed orders as the fixed tolerance.

\subsection{Local-to-global transfer and main theorem}

\begin{lemma}[Local-to-global transfer]
\label{lem:local-global-transfer}
Let $\Psi_h$ be a one-step map on the reduced chart, $t_n=nh$, and $y_{n+1}=\Psi_h(y_n)$ with $y_0$
the exact initial state. Suppose that for $0\le n<N$ with $Nh\le T$ the exact states
$Y_n=\varphi_{t_n}(y_0)$ satisfy $\|\Psi_h(Y_n)-Y_{n+1}\|\le C_\ell h^{p+1}$, that
$\|\Psi_h(y_n)-\Psi_h(Y_n)\|\le(1+C_sh)\|y_n-Y_n\|$ whenever $y_n\in K$, and that
$C_\ell\Gamma_s(T)h^p\le d_K$, where $\Gamma_s(T)=(e^{C_sT}-1)/C_s$ for $C_s>0$ and
$\Gamma_s(T)=T$ for $C_s=0$. Then $y_n\in K$ for all $n\le N$ and
\begin{equation}
\label{eq:local-global-lemma-grid-bound}
  \max_{0\le n\le N}\|y_n-Y_n\|\le C_\ell\Gamma_s(T)h^p .
\end{equation}
If moreover the reporting map $\mathcal R$ (chart to reported position and velocity) is
$C_{\mathcal R}$-Lipschitz on $K$, then
$\max_n\|\mathcal R(y_n)-\mathcal R(Y_n)\|\le C_{\mathcal R}C_\ell\Gamma_s(T)h^p$.
\end{lemma}

\begin{proof}
Let $e_n=\|y_n-Y_n\|$. By induction, as long as $y_n\in K$,
$e_{n+1}\le(1+C_sh)e_n+C_\ell h^{p+1}$, hence $e_n\le C_\ell h^{p+1}\sum_{k<n}(1+C_sh)^k$. For
$C_s>0$ the geometric sum equals $((1+C_sh)^n-1)/(C_sh)\le(e^{C_snh}-1)/(C_sh)\le\Gamma_s(T)/h$,
and for $C_s=0$ it equals $n\le T/h$. Thus $e_n\le C_\ell\Gamma_s(T)h^p\le d_K$, which keeps $y_n$
in the tube and closes the induction. The reporting bound is the Lipschitz property. The lemma is
machine-checked.
\end{proof}

\begin{definition}[Branch-selected one-step map]
\label{def:accepted-branch-map}
The one-step map $\Psi_h$ is defined on the set $\mathcal D_h\subset K$ of reduced states for which
the transition satisfies (H2) and (H3): the computed iterate lies in $B_\rho(Z_G)$ with residual at
most $\eta_h$, the endpoint is reconstructed by Eqs.~\eqref{eq:end-trans}--\eqref{eq:end-rot} and
closed by $\mathcal C_h$, and $\Psi_h(y)=\mathcal P_h(\widetilde Z)$.
\end{definition}

\begin{theorem}[Sixth-order error bound]
\label{thm:g6fullva-order}
Let Assumption~\ref{ass:regularity} hold and let $C_{\rm loc}=C_G+C_E+C_Nc_\eta$ with $C_G$, $C_E$,
$C_N$ from Lemmas~\ref{lem:gauss-endpoint-defect}, \ref{lem:endpoint-closure}
and~\ref{lem:inexact-newton}. Then there is $h_0>0$ such that for $0<h\le h_0$ and every
$y\in\mathcal D_h$,
\begin{equation}
\label{eq:g6fva-local-defect-theorem}
  \|\Psi_h(y)-\varphi_h(y)\|\le C_{\rm loc}h^7 .
\end{equation}
Consequently, for every grid $t_n=nh$, $0\le n\le N_h$, $N_hh\le T$, whose transitions lie in
$\mathcal D_h$ and whose numerical trajectory $y_{n+1}=\Psi_h(y_n)$ starts from the exact reduced
initial state,
\begin{equation}
\label{eq:g6fva-reduced-grid-bound}
  \max_{0\le n\le N_h}\|y_n-y(t_n)\|\le C_{\rm red}h^6,\qquad C_{\rm red}=C_{\rm loc}\Gamma_s(T),
\end{equation}
and the reported position--velocity map $\mathcal R(y)=(q(y),v(y))$ satisfies
\begin{equation}
\label{eq:g6fva-reported-grid-bound}
  \max_{0\le n\le N_h}\|\mathcal R(y_n)-\mathcal R(y(t_n))\|\le C_{qv}h^6,\qquad
  C_{qv}=C_{\mathcal R}C_{\rm red} .
\end{equation}
The constants depend on the tube, on $T$ and on the constants of Assumption~\ref{ass:regularity},
including $c_\eta$, but not on $h$, $N_h$ or the linear-algebra backend. Under (H1) alone the same
conclusions hold for $h\le h_0$ whenever the stage solve is started in $B_{\rho/2}(Z_G)$ and
stopped by the rule $\|F_h\|\le c_\eta h^7$ (Lemmas~\ref{lem:p2-from-p1}
and~\ref{lem:newton-envelope}).
\end{theorem}

\begin{proof}
Fix $y\in\mathcal D_h$ and write $\Psi_h^G=\mathcal E_h(Z_G)$, $\Psi_h^E=\mathcal P_h(Z_G)$ and
$\Psi_h(y)=\mathcal P_h(\widetilde Z)$. By Lemma~\ref{lem:stage-residual-defect} the exact stage
root on the branch is $Z_G$, so there is no stage-root perturbation term, and
\begin{equation}
\label{eq:three-term-local-defect-sum}
\begin{aligned}
  \|\Psi_h(y)-\varphi_h(y)\|
  &\le\|\Psi_h^G(y)-\varphi_h(y)\|+\|\Psi_h^E(y)-\Psi_h^G(y)\|\\
  &\quad+\|\Psi_h(y)-\Psi_h^E(y)\|
  \le C_Gh^7+C_Eh^7+C_Nc_\eta h^7
\end{aligned}
\end{equation}
by Lemmas~\ref{lem:gauss-endpoint-defect}, \ref{lem:endpoint-closure}
and~\ref{lem:inexact-newton}, after shrinking $h_0$ so that all three apply on the tube. This is
Eq.~\eqref{eq:g6fva-local-defect-theorem}. Choose $h_0$ additionally so that
$C_{\rm loc}\Gamma_s(T)h_0^6\le d_K$. Lemma~\ref{lem:local-global-transfer} with $p=6$,
$C_\ell=C_{\rm loc}$ and the stability scale Eq.~\eqref{eq:assumption-stability-scale} gives
Eq.~\eqref{eq:g6fva-reduced-grid-bound}; the reporting-map clause gives
Eq.~\eqref{eq:g6fva-reported-grid-bound}. The last sentence follows from
Lemmas~\ref{lem:p2-from-p1} and~\ref{lem:newton-envelope}, which supply (H2) and (H3) from (H1) for
small steps.
\end{proof}

Two remarks on the scope of the theorem. First, it is an estimate for the reduced state and for the
reported positions and velocities; it says nothing about the accuracy of the stage multipliers or
reaction forces beyond what follows from the smooth dependence of the lift on the reduced state, and
it does not turn a small constraint residual on a closed-loop mechanism into a trajectory-error
bound, which would need a stability constant for that mechanism. Second, because the stage residual
vanishes identically at the lifted Gauss stage, the order of the method is that of reduced Gauss
collocation; the absolute-coordinate formulation contributes the exact three-level constraint
enforcement, the multipliers and reaction forces at every stage, the friction law through the
multipliers, and the Lie-group endpoint reconstruction, but no additional stage-level
approximation. The same argument gives order $2s$ for the $s$-stage Gauss member with the same
stage system, which Section~\ref{sec:numerical} confirms for $s=1,2$.

\paragraph{Formal-order comparison with time finite elements}
The Lie-group time finite-element family of~\cite{chaturvedi2026tfe} with Gauss--Lobatto
polynomials of degree $m$ has formal order $2m-1$~\cite{kim2017higher}; its $m=3$ member is a
fifth-order formula. The comparison of the present method with that family is at this formal level
only: order six against order five for a stage system of comparable size. We have not implemented
the time finite-element method and make no statement about its accuracy or cost on its own test
problems.
